%% kapa: exam=1 round=33 year=2022 period=2 form=A subject=appraisal_law subject_ko=감정평가관계법규 seq=21 number=21
%% source: Q-Net 감정평가사 기출문제 (cst00309) artlSeq=5209997 file=raw/1008/5209997/2022년도 제33회 감정평가사 1차시험 2교시(감정평가 관계법규, 회계학).pdf
%% conversion: kordoc 4.24.51 -> tools/gen1.py (auto); review=proofread-sonnet; flags=none
\begin{question}{21}
국유재산법령상 행정재산에 관한 설명으로 옳은 것은?
\begin{choices}
  \item 중앙관서의 장은 사용허가한 행정재산을 지방자치단체가 직접 공용으로 사용하기 위하여 필요하게 된 경우에도 그 허가를 철회할 수 없다.
  \item 행정재산의 관리위탁을 받은 자가 그 재산의 일부를 사용·수익하는 경우에는 미리 해당 중앙관서의 장의 승인을 받아야 한다.
  \item 경작용으로 실경작자에게 행정재산의 사용허가를 하려는 경우에는 일반경쟁에 부쳐야 한다.
  \item 수의의 방법으로 한 사용허가는 허가기간이 끝난 후 갱신할 수 없다.
  \item 행정재산의 사용허가를 한 날부터 3년 내에는 사용료를 조정할 수 없다.
\end{choices}
\end{question}
\answer{2}
